\usection{Análisis de la Problemática}

Esta fase inicial se centró en una indagación profunda para establecer las bases teóricas y técnicas del proyecto. Se llevó a cabo una revisión documental que abarcó desde los fundamentos del pensamiento computacional hasta las arquitecturas de software necesarias para soportar la ejecución de flujos de datos en entornos web y distribuidos, manteniendo un balance entre capacidad técnica y usabilidad.

Se analizaron los marcos de trabajo existentes para la enseñanza de la computación, tomando como referencia principal la propuesta de \citeA{brennan2012} en el lenguaje de programación Scratch. Se identificaron los conceptos clave que el lenguaje debe ser capaz de modelar, tales como secuencias, bucles, eventos, condicionales, operadores y datos. Se puso especial énfasis en el concepto de concurrencia, pudiendo identificarlo como una carencia en las herramientas educativas tradicionales que suelen enfocarse en la ejecución secuencial estricta. Asimismo, se consideraron las definiciones de \citeA{zapataros2015} sobre la importancia de la abstracción, el reconocimiento de patrones y la creación de algoritmos como pilares para la resolución de problemas.

Por otro lado, para asegurar que la herramienta fuese adecuada para la población objetivo (niños y jóvenes de 8 a 16 años), se revisaron las teorías del desarrollo cognitivo de Jean Piaget. Se determinó que el diseño visual debe apoyar el paso de la etapa de Operaciones Concretas hasta las Operaciones Formales, permitiendo a los estudiantes manipular representaciones tangibles de datos para comprender conceptos abstractos. En la tabla \ref{tab:analisis-poblacion} se resumen las características cognitivas de esta población y cómo el diseño del lenguaje debe adaptarse a ellas.

\begin{table}[H]
  \centering
  \caption{Caracterización de la población objetivo según las etapas de desarrollo cognitivo de Piaget}
  \label{tab:analisis-poblacion}
  \begingroup
  \renewcommand{\arraystretch}{\TableRowStretch}
  \renewcommand{\TableBodyFontSize}{\footnotesize}
  \begin{tabularx}{\textwidth}{ X X X }
    \hline
    \textbf{Aspecto de Diseño}                          & \textbf{Niños (8-11 años)}                                                                                  & \textbf{Jóvenes (12-16 años)}                                                                   \\
    \hline
    \textbf{Fase cognitiva según Piaget}                & Operaciones Concretas                                                                                       & Operaciones Formales                                                                            \\
    \textbf{Características cognitivas}                 & Pensamiento lógico anclado en la experiencia real y lo observable. Dificultad con razonamientos abstractos. & Razonamiento lógico en todas las circunstancias. Capacidad de manejar lo simbólico y abstracto. \\
    \textbf{Orientación de la interfaz}                 & Facilidad de entrada para principiantes, eliminando barreras de inicio (Piso Bajo).                         & Capacidad de gestionar complejidad y permitir una progresión avanzada (Techo Alto).             \\
    \textbf{Manejo de abstracción}                      & Uso de metáforas visuales claras y limitación de nodos disponibles para no abrumar.                         & Soporte para la modularidad y encapsulación de flujos.                                          \\
    \textbf{Retroalimentación y visualización de datos} & Feedback inmediato, tipificación visual de datos para facilitar la comprensión.                             & Herramientas de depuración de alto nivel, trazabilidad e inspección de estado.                  \\
    \textbf{Lógica y control de flujo}                  & Refuerzo visual de la causalidad y el flujo de datos.                                                       & Soporte para lógica condicional y estructuras de control avanzadas.                             \\
    \hline
  \end{tabularx}
  \endgroup
\end{table}

Adicionalmente, se estudió la teoría del Construccionismo de Seymour Papert, específicamente su visión sobre los “micromundos” y la “Sociedad de la Mente”. Se rescató la idea de que los sistemas complejos pueden entenderse mejor como una colección de agentes distribuidos que interactúan entre sí a través del paso de mensajes, lo cual fundamenta la necesidad de soportar concurrencia y sistemas distribuidos en la propuesta.

\usubsection{Definición de requerimientos}

Producto del análisis del problema, el cual incluyó la revisión de la literatura relacionada y la entrevista con el profesor Jesús Lárez, especialista en desarrollo del pensamiento computacional y lenguajes de programación, se definieron los siguientes requerimientos:

\usubsubsection{Requerimientos funcionales}

\begin{itemize}
  \item El lenguaje de programación debe ser Turing completo.
  \item El lenguaje debe ofrecer soporte explícito de diferentes tipos de datos, tanto simples como compuestos, con conversión automática entre ellos.
  \item El lenguaje debe representar visualmente los tipos de datos (números, booleanos, cadenas, listas) mediante el uso de estilos distintos en las conexiones y nodos.
  \item El lenguaje debe permitir abstracciones funcionales y modularidad, agrupando conjuntos de nodos para encapsularlos en subflujos.
  \item El entorno de desarrollo y ejecución debe ser interactivo y debe mostrar retroalimentación visual continua sobre el estado de cada nodo (activo, ejecutado, esperando datos, o en error).
  \item El entorno debe mostrar en tiempo real el movimiento y transformación de los datos a través de las conexiones.
  \item El entorno de ejecución debe facilitar la depuración de los programas realizados.
  \item El entorno debe soportar \emph{literatura}, permitiendo compartir programas con otros usuarios.
  \item El lenguaje debe ejecutar los nodos activando operaciones automáticamente cuando sus datos de entrada estén disponibles.
  \item El entorno debe permitir la ejecución inmediata del programa.
  \item El entorno debe permitir la exportación e importación de proyectos para facilitar que los usuarios construyan sobre el código de otros.
  \item El entorno debe ofrecer un módulo de desafíos de programación estructurados cortos e interactivos, centrados en el uso de un conjunto específico de nodos, con el objetivo de fomentar el aprendizaje significativo de forma lúdica.
  \item El entorno debe promover el interés por el lenguaje via ramificación de desafíos de acuerdo a la edad del usuario.
  \item El entorno y el lenguaje deben soportar la programación distribuida, permitiendo ejecución remota de partes de código y funcionalidades nativas asociadas.
  \item El entorno debe ofrecer de forma nativa soporte de computación física, de entrada para tarjetas micro:bit.
  \item El entorno debe ofrecer las construcciones del lenguaje categorizadas por diferentes niveles de abstracción.
\end{itemize}

\usubsubsection{Requerimientos no funcionales}

\begin{itemize}
  \item Se debe ofrecer una interfaz que sea atractiva e intuitiva para niños y jóvenes de 8 a 16 años.
  \item Debe ser una herramienta concebida bajo los principios de piso bajo, techo alto y muros anchos.
  \item El entorno se debe implementar como una aplicación basada en tecnologías web.
  \item Los nodos deben ser simples y de fácil lectura, indicando su aridad, tipos de datos de entrada y salida, y la función que realizan.
  \item Los nodos deben estar categorizados y codificados en diferentes colores, con descripción de su funcionalidad.
  \item El lenguaje debe ser consistente en la implementación de sus construcciones, evitando casos especiales que no se apliquen de forma general.
  \item El lenguaje debe ser extensible, permitiendo la adición de nuevos nodos y funcionalidades sin necesidad de modificar la arquitectura central del sistema.
  \item El entorno y el lenguaje deben fomentar la experimentación y el aprendizaje iterativo.
  \item La interfaz debe mantener una correspondencia entre la representación visual y la lógica de ejecución, asegurando que no existan estados ocultos al usuario.
  \item Se debe fomentar su uso en un contexto social, permitiendo interacciones que apoyen el aprendizaje colaborativo.
\end{itemize}
